% ============================================================
% plan_ahead.tex
% Plan-Ahead Cluster Optimization Model — LaTeX Reference
%
% Generated from plan_ahead_source_of_truth.txt
% Compile with: pdflatex plan_ahead.tex  (run twice for TOC)
% ============================================================

\documentclass[11pt,a4paper]{article}

% ── Packages ─────────────────────────────────────────────────
\usepackage[T1]{fontenc}
\usepackage[utf8]{inputenc}
\usepackage{lmodern}
\usepackage[margin=2.5cm]{geometry}
\usepackage{amsmath,amssymb,amsthm}
\usepackage{mathtools}        % \coloneqq, dcases
\usepackage{booktabs}         % \toprule, \midrule, \bottomrule
\usepackage{array}            % column types in tabular
\usepackage{longtable}        % tables spanning multiple pages
\usepackage{xcolor}
\usepackage{hyperref}
\usepackage{enumitem}
\usepackage{microtype}
\usepackage{parskip}          % space between paragraphs, no indent
\usepackage{titlesec}

% ── Hyperref setup ────────────────────────────────────────────
\hypersetup{
  colorlinks = true,
  linkcolor  = blue!60!black,
  urlcolor   = blue!60!black,
  citecolor  = green!50!black,
  pdftitle   = {Plan-Ahead Cluster Optimization Model},
  pdfauthor  = {Alrick Grandison},
}

% ── Custom colours ────────────────────────────────────────────
\definecolor{constblue}{RGB}{30,80,160}
\definecolor{notegrey}{RGB}{90,90,90}
\definecolor{fixred}{RGB}{180,30,30}

% ── Custom commands ───────────────────────────────────────────
\newcommand{\cT}{\mathcal{T}}           % tenants
\newcommand{\cW}{\mathcal{W}}           % workloads
\newcommand{\cN}{\mathcal{N}}           % nodes
\newcommand{\cR}{\mathcal{R}}           % resources
\newcommand{\cH}{\mathcal{H}}           % time slots
\newcommand{\cQ}{\mathcal{Q}}           % QoS classes
\newcommand{\cK}{\mathcal{K}}           % isolation primitives

\newcommand{\kap}{\kappa}               % kappa (Cantelli)
\newcommand{\sig}{\sigma}               % sigma (fairness)
\newcommand{\eps}{\varepsilon}          % epsilon (risk)
\newcommand{\BigM}{\mathsf{M}}          % big-M constant

% Constraint label box
\newcommand{\clabel}[1]{%
  \textcolor{constblue}{\textbf{(#1)}}}

% Note environment
\newenvironment{note}{\smallskip\noindent\textcolor{notegrey}{\textit{Note:\ }}}{\smallskip}

% Implementation note
\newenvironment{impl}{\smallskip\noindent\textcolor{fixred}{\textit{Implementation note:\ }}}{\smallskip}

% ── Section style ─────────────────────────────────────────────
\titleformat{\section}{\large\bfseries}{\thesection.}{0.5em}{}[\vspace{-0.3em}\rule{\linewidth}{0.4pt}]
\titleformat{\subsection}{\normalsize\bfseries}{\thesubsection.}{0.5em}{}

% ─────────────────────────────────────────────────────────────
\title{\textbf{Plan-Ahead Cluster Optimization Model}\\[0.4em]
  \large Mixed-Integer Second-Order Cone Program (MISOCP)}
\author{DAMO 699 Capstone $\cdot$ Group 3\\
        Multi-Tenant Cluster Memory Scheduling}
\date{\today}
% ─────────────────────────────────────────────────────────────

\begin{document}

\maketitle
\tableofcontents
\newpage

% ============================================================
\section{Introduction}
% ============================================================

The plan-ahead model solves a periodic resource reservation problem for a
multi-tenant compute cluster.  It is invoked once per planning horizon
$\cH$ (e.g.\ once per day, once per week, or once every three days --- the
horizon length is configurable) and produces a schedule that assigns each
admitted tenant a set of authorised cluster nodes for each time slot.

The real-time optimizer is invoked every scheduling epoch (e.g.\ every
60~seconds) to place individual arriving jobs.  It receives the authorised
node sets $A_t^{(i)}$ from the plan-ahead output and enforces them via
constraint~\ref{c5}.

\paragraph{System pipeline.}
\[
  \underbrace{\text{Prediction}}_{\mu_{i,j,r},\,\Sigma_r,\,d_{i,j,r}}
  \;\longrightarrow\;
  \underbrace{\text{Plan-Ahead (MISOCP)}}_{\text{Gurobi, once per }\cH}
  \;\longrightarrow\;
  A_t^{(i)}
  \;\longrightarrow\;
  \underbrace{\text{Real-Time (MILP)}}_{\text{OR-Tools, every epoch}}
\]

The model simultaneously decides:
\begin{itemize}[nosep]
  \item which tenants to admit ($a_i$);
  \item where to place each admitted workload ($x_{i,j,n,t}$);
  \item which isolation primitive each workload uses ($w_{i,j,k,t}$);
  \item how many nodes to activate ($z_{n,t}$);
  \item when workloads migrate between time slots ($m_{i,j,n,t}^{\text{mig}}$).
\end{itemize}

% ============================================================
\section{Sets}
% ============================================================

\begin{longtable}{@{}lll@{}}
  \toprule
  \textbf{Symbol} & \textbf{Definition} & \textbf{Description} \\
  \midrule
  \endhead
  $\cT$ & $\{0,1,\ldots,|\cT|-1\}$ & Tenants sharing the cluster \\
  $\cW_i$ & $\{0,1,\ldots,|\cW_i|-1\}$ & Workloads of tenant $i$ \\
  $\cN$ & $\{0,1,\ldots,|\cN|-1\}$ & Cluster nodes \\
  $\cR$ & $\{0\text{=CPU},\;1\text{=MEM}\}$ & Resource types \\
  $\cH$ & $\{0,1,\ldots,|\cH|-1\}$ & Planning time slots \\
  $\cQ$ & $\{0\text{=guaranteed},\;1\text{=burstable}\}$ & QoS classes \\
  $\cK$ & $\{0\text{=none},\;1\text{=gVisor},\;2\text{=Kata}\}$ & Isolation primitives \\[4pt]
  $\mathcal{P}$ & $\{(i,j,i',j') : (i,j)<(i',j')\}$ & Ordered pairs of distinct workloads \\
  & & (avoids duplicate co-location variables) \\
  $N_{\text{allowed}}^{(i,j)}$ & $\subseteq \cN$ & Nodes where workload $(i,j)$ may legally run \\
  & & (data-residency / compliance constraint) \\
  \bottomrule
\end{longtable}

% ============================================================
\section{Decision Variables}
% ============================================================

\begin{longtable}{@{}p{4.2cm}p{2.1cm}p{8.0cm}@{}}
  \toprule
  \textbf{Variable} & \textbf{Domain} & \textbf{Description} \\
  \midrule
  \endhead
  $x_{i,j,n,t}$ & $\{0,1\}$
    & 1 if workload $j$ of tenant $i$ is placed on node $n$ at slot $t$ \\[3pt]
  $y_{i,n,t}$ & $\{0,1\}$
    & 1 if tenant $i$ has at least one workload on node $n$ at slot $t$ \\[3pt]
  $z_{n,t}$ & $\{0,1\}$
    & 1 if node $n$ is active (hosting any workload) at slot $t$ \\[3pt]
  $w_{i,j,k,t}$ & $\{0,1\}$
    & 1 if isolation primitive $k$ is used for workload $(i,j)$ at slot $t$ \\[3pt]
  $m_{i,j,n,t}^{\text{mig}}$ & $\{0,1\}$
    & 1 if workload $(i,j)$ migrates \emph{onto} node $n$ at slot $t$
      ($t\ge 1$ only; forced~$0$ at $t=0$) \\[3pt]
  $a_i$ & $\{0,1\}$
    & 1 if tenant $i$ is admitted for the planning horizon \\[3pt]
  $e_{i,j,t}$ & $\mathbb{R}_{\ge 0}$
    & Excess latency beyond SLA target $L_{i,q(i,j)}$ for workload $(i,j)$ at $t$ \\[3pt]
  $\eps_{n,t}^{\text{eff}}$ & $[0,1]$
    & Effective SLA risk level on node $n$ at $t$
      (driven to $\min_i \eps_i$ among tenants present) \\[3pt]
  $\xi_{i,j,n,k,t}$ & $\{0,1\}$
    & McCormick auxiliary: $\xi_{i,j,n,k,t} = x_{i,j,n,t}\cdot w_{i,j,k,t}$ \\[3pt]
  $\zeta_{i,j,i',j',n,t}$ & $\{0,1\}$
    & McCormick auxiliary: $\zeta = x_{i,j,n,t}\cdot x_{i',j',n,t}$
      (co-location on same node); defined for $(i,j)<(i',j')$ \\[3pt]
  $\sigma$ & $[0,1]$
    & Fairness variable: min DRF satisfaction across admitted tenants \\
  \bottomrule
\end{longtable}

% ============================================================
\section{Parameters}
% ============================================================

\subsection{Resource Demand}

\begin{longtable}{@{}p{2.2cm}p{11.5cm}@{}}
  \toprule
  \textbf{Symbol} & \textbf{Description} \\
  \midrule
  \endhead
  $d_{i,j,r}$   & Declared resource demand of workload $(i,j)$ for resource $r$.
                   Used as the numerator in DRF satisfaction (\S\ref{sec:c7}). \\[4pt]
  $\mu_{i,j,r}$ & Mean predicted resource \emph{usage} of workload $(i,j)$ for
                   resource $r$.  Estimated from historical traces
                   (\texttt{InstanceUsage.average\_usage}).
                   Used in the probabilistic capacity mean-load term (\S\ref{sec:socp}). \\[4pt]
  $\Sigma_r$    & Empirical covariance matrix of resource $r$ across all workloads,
                   shape $|\cW_{\text{total}}|\times|\cW_{\text{total}}|$, positive
                   semi-definite.  Estimated from stochastic usage samples.
                   Used in the SOCP safety-buffer term. \\[4pt]
  $C_{n,r}$     & Capacity of node $n$ for resource $r$
                   (\texttt{MachineEvents.capacity}). \\[4pt]
  $Q_{i,r}$     & Resource quota for tenant $i$, resource $r$
                   (denominator in DRF satisfaction).
                   Set to $C_{\cdot,r}/2$ in the current implementation. \\[4pt]
  $\alpha_r^{\text{oc}}$ & Overcommit ratio for resource $r$.  Available for
                   soft-overcommit extensions; currently only $C_{n,r}$ is used
                   as the RHS capacity. \\[4pt]
  $N_{i,t}$      & Expected number of active jobs (workloads) for tenant $i$
                   during planning slot $t$.  Used to size $|\cW_i|$ before the
                   model is built: $|\cW_i| = N_{i,t}$ (or the maximum across
                   $t\in\cH$ if a single set is used for all slots).
                   Derived from trace by counting \texttt{SUBMIT} events per
                   tenant per 4-hour bucket.  Autopilot-managed workloads are
                   excluded (their task count is cluster-controlled). \\
  \bottomrule
\end{longtable}

\subsection{Pricing and Risk}

\begin{longtable}{@{}p{2.2cm}p{11.5cm}@{}}
  \toprule
  \textbf{Symbol} & \textbf{Description} \\
  \midrule
  \endhead
  $p_{i,j}$           & SLA penalty weight for workload $(i,j)$.  Coefficient on
                         $e_{i,j,t}$ in the objective. \\[4pt]
  $\pi_n$             & Infrastructure cost per active node-slot.  Coefficient on
                         $z_{n,t}$ in the objective. \\[4pt]
  $\bar{\pi}_i$       & Gross revenue from admitting tenant $i$ (billing contract value).
                         Scales with planning horizon: $\bar{\pi}_i = 3|\cH|$
                         in the current implementation. \\[4pt]
  $v_i^{\text{op}}$   & Operational cost of serving tenant $i$ (support overhead).
                         Net revenue per tenant: $\bar{\pi}_i - v_i^{\text{op}}$.
                         Scales with horizon: $v_i^{\text{op}} = 0.5|\cH|$. \\[4pt]
  $\eps_i$            & SLA risk tolerance for tenant $i$ --- maximum allowable
                         probability of capacity overrun used in the Cantelli bound.
                         Smaller $\eps_i$ is stricter (larger safety margin). \\
  \bottomrule
\end{longtable}

\subsection{Isolation Primitives}

Isolation primitives model Kubernetes security isolation mechanisms:
$k=0$~(none: Linux namespaces), $k=1$~(gVisor: user-space kernel),
$k=2$~(Kata: VM-based containers).

\begin{longtable}{@{}p{2.2cm}p{11.5cm}@{}}
  \toprule
  \textbf{Symbol} & \textbf{Description} \\
  \midrule
  \endhead
  $\eta_{k,r}$        & Resource overhead multiplier of primitive $k$ for resource
                         $r$.  Effective usage $= \eta_{k,r}\,\mu_{i,j,r}\,\xi_{i,j,n,k,t}$.
                         Values: $\eta_{0,r}=1.0$, $\eta_{1,r}=1.20$ (gVisor $+20\%$),
                         $\eta_{2,r}=1.05$ (Kata $+5\%$). \\[4pt]
  $c_k$               & Operational cost premium for running primitive $k$. \\[4pt]
  $\rho_{k,k'}$       & Residual interference between workloads using primitives $k$
                         and $k'$ on the same node, $\rho\in[0,1]$.
                         $\rho=0$: fully isolated; $\rho=1$: no isolation. \\[4pt]
  $\tau_{i,i'}$       & Maximum acceptable residual interference between tenants
                         $i$ and $i'$.  Constraint C3c enforces
                         $\rho_{k,k'}\le\tau_{i,i'}$ for co-located workloads. \\[4pt]
  $k_{\min,i}$        & Minimum isolation primitive level required for tenant $i$
                         (compliance floor, e.g.\ regulatory mandate for gVisor). \\
  \bottomrule
\end{longtable}

\subsection{Control-Plane Budgets}

Control-plane resources (etcd writes, API-server QPS, service objects) are
shared cluster-wide and capped per time slot.

\begin{longtable}{@{}p{3.2cm}p{10.5cm}@{}}
  \toprule
  \textbf{Symbol} & \textbf{Description} \\
  \midrule
  \endhead
  $\omega_{i,k}^{\text{etcd}}$ & etcd write operations consumed per workload of
                                   tenant $i$ with primitive $k$. \\[4pt]
  $\omega_{i,k}^{\text{qps}}$  & API-server QPS consumed per workload of tenant $i$
                                   with primitive $k$. \\[4pt]
  $\omega_{i,k}^{\text{svc}}$  & Kubernetes service objects consumed per workload of
                                   tenant $i$ with primitive $k$. \\[4pt]
  $B^{\text{etcd}}$            & Cluster-wide etcd write budget per time slot. \\[4pt]
  $B^{\text{qps}}$             & Cluster-wide API-server QPS budget per time slot. \\[4pt]
  $B^{\text{svc}}$             & Cluster-wide service object budget per time slot. \\[4pt]
  $\delta^{\text{qps}}$        & QPS churn cost per migration event (endpoint updates). \\[4pt]
  $B^{\text{qps,churn}}$       & Maximum allowed migration-induced QPS churn per slot. \\
  \bottomrule
\end{longtable}

\subsection{SLA and Migration}

\begin{longtable}{@{}p{2.2cm}p{11.5cm}@{}}
  \toprule
  \textbf{Symbol} & \textbf{Description} \\
  \midrule
  \endhead
  $q(i,j)$      & QoS class assigned to workload $(i,j)$.
                   Determines the latency target $L_{i,q}$. \\[4pt]
  $L_{i,q}$     & Latency SLA target for tenant $i$, QoS class $q$ (milliseconds).
                   C5 enforces modelled latency $\le L_{i,q}+e_{i,j,t}$. \\[4pt]
  $\gamma_{i,j}$& Migration cost weight for workload $(i,j)$.  Multiplied by
                   $m_{i,j,n,t}^{\text{mig}}$ in the objective. \\[4pt]
  $\Delta_i$    & Per-tenant disruption budget --- maximum workload migrations
                   allowed per time slot. \\
  \bottomrule
\end{longtable}

\subsection{Latency Model}

\begin{longtable}{@{}p{3.0cm}p{10.7cm}@{}}
  \toprule
  \textbf{Symbol} & \textbf{Description} \\
  \midrule
  \endhead
  $l^0_{i,j}$                        & Base latency of workload $(i,j)$ running
                                        in isolation (ms). \\[4pt]
  $\alpha_{i,j,i',j'}^{\text{lat}}$  & Latency increase on $(i,j)$ due to
                                        co-location with workload $(i',j')$ on the
                                        same node.  Modelled linearly via $\zeta$. \\[4pt]
  $\beta_{i,j,k}^{\text{lat}}$       & Latency overhead added by isolation primitive
                                        $k$ for workload $(i,j)$ (e.g.\ gVisor syscall
                                        overhead). \\
  \bottomrule
\end{longtable}

\subsection{Objective Weights and SOCP Auxiliaries}

\begin{longtable}{@{}p{2.6cm}p{11.1cm}@{}}
  \toprule
  \textbf{Symbol} & \textbf{Description} \\
  \midrule
  \endhead
  $\lambda_s,\;s\in\{0,\ldots,5\}$
    & Objective weights.  Defaults: $\lambda_0=1.0$ (infra), $\lambda_1=10.0$
      (SLA), $\lambda_2=1.0$ (revenue), $\lambda_3=5.0$ (fairness),
      $\lambda_4=0.5$ (isolation), $\lambda_5=2.0$ (migration). \\[4pt]
  $\kap(\eps)$
    & Cantelli one-sided quantile factor:
      $\kap(\eps)=\sqrt{(1-\eps)/\eps}$.
      $\kap(0.05)\approx 4.36$ (95\% confidence),
      $\kap(0.10)\approx 3.00$ (90\% confidence). \\[4pt]
  $\kap_n$
    & Per-node kappa (implementation):
      $\kap_n = \kap\!\left(\min\{\eps_i : n\in N_{\text{allowed}}^{(i,j)}
      \text{ for some }j\}\right)$.
      Less conservative than the global minimum. \\[4pt]
  $L_r$
    & Lower-triangular Cholesky factor of $\Sigma_r$:
      $\Sigma_r = L_r L_r^\top$. \\[4pt]
  $r^*_i$
    & Dominant resource for tenant $i$ (DRF):
      $r^*_i = \arg\max_r \sum_j d_{i,j,r}/Q_{i,r}$. \\[4pt]
  $s_i$
    & DRF satisfaction fraction for tenant $i$:
      $s_i = \frac{\sum_{j,n} d_{i,j,r^*_i}\,x_{i,j,n,t}}{Q_{i,r^*_i}}$. \\[4pt]
  $A_t^{(i)}$
    & Authorised node set for tenant $i$ at time $t$ (model output):
      $A_t^{(i)} = \{n\in\cN : \sum_j x_{i,j,n,t} \ge 1\}$. \\
  \bottomrule
\end{longtable}

% ============================================================
\section{Probabilistic Capacity Constraint (SOCP)}
\label{sec:socp}
% ============================================================

Because workload resource usage is stochastic, a deterministic safety buffer
is derived via the \textbf{Cantelli inequality} (one-sided Chebyshev).  For
any random variable $X$ with mean $\mu$ and variance $\sigma^2$:
\[
  \Pr\!\bigl[X > \mu + \kap(\eps)\,\sigma\bigr] \;\le\; \eps,
  \qquad
  \kap(\eps) = \sqrt{\frac{1-\eps}{\eps}}.
\]

For workloads placed on node $n$, the joint resource usage is a sum of
correlated random variables.  The multivariate version uses the covariance
matrix $\Sigma_r$.  After Cholesky decomposition $\Sigma_r = L_r L_r^\top$,
the safety-buffer term becomes an $\ell_2$ norm --- a \textbf{second-order
cone} constraint, making the full model a MISOCP.

\paragraph{Constraint C2.}\label{c2}
For each node $n\in\cN$, resource $r\in\cR$, time slot $t\in\cH$:
\begin{align}
  \underbrace{\sum_{i,j,k} \eta_{k,r}\,\mu_{i,j,r}\,\xi_{i,j,n,k,t}}_{\text{mean load}}
  \;+\;
  \kap_n \underbrace{\bigl\|L_r\,\boldsymbol{\xi}_{n,t}\bigr\|_2}_{\text{safety buffer}}
  &\;\le\; C_{n,r}\,z_{n,t},
  \label{eq:c2}
\end{align}
where $\boldsymbol{\xi}_{n,t}\in\{0,1\}^{|\cW_{\text{total}}|}$ has entry
$q$ equal to $\sum_k \xi_{\text{wl}(q),n,k,t}$ (total activation of
workload $q$ on node $n$ at $t$).

The $\ell_2$ cone is introduced via auxiliary variables:
\begin{align*}
  v_s &= \sum_{q} [L_r]_{sq}\,(\boldsymbol{\xi}_{n,t})_q,
  \quad s=1,\ldots,|\cW_{\text{total}}|, \\
  v_s^2 &\le \phi^2 \quad(\text{i.e.,}\;\|\mathbf{v}\|_2 \le \phi),
\end{align*}
then replace the $\ell_2$ term with $\kap_n\,\phi$. Gurobi encodes this
directly as a quadratic constraint (SOCP).

\paragraph{Per-node effective risk.}
The auxiliary $\eps_{n,t}^{\text{eff}}$ tracks the strictest tenant on each node:
\[
  \eps_{n,t}^{\text{eff}} \;\le\; \eps_i + (1 - y_{i,n,t})
  \quad \forall\,i\in\cT.
\]
This drives $\eps_{n,t}^{\text{eff}}$ to $\min_i\{\eps_i : y_{i,n,t}=1\}$.
In the current implementation $\kap_n$ is precomputed (no binary dependency)
using the set $N_{\text{allowed}}^{(i,j)}$ as a conservative proxy.

% ============================================================
\section{McCormick Linearizations}
% ============================================================

Two bilinear products of binary variables appear; both are linearized via
standard McCormick envelopes.

\subsection{$\xi_{i,j,n,k,t} = x_{i,j,n,t}\cdot w_{i,j,k,t}$}
Used in: mean load (C2), latency primitive overhead (C5).
\begin{align*}
  \xi_{i,j,n,k,t} &\le x_{i,j,n,t}, \\
  \xi_{i,j,n,k,t} &\le w_{i,j,k,t}, \\
  \xi_{i,j,n,k,t} &\ge x_{i,j,n,t} + w_{i,j,k,t} - 1.
\end{align*}

\subsection{$\zeta_{i,j,i',j',n,t} = x_{i,j,n,t}\cdot x_{i',j',n,t}$}
Defined for ordered pairs $(i,j)<(i',j')$ to avoid duplication.
Used in: latency interference (C5), isolation compatibility (C3c).
\begin{align*}
  \zeta_{i,j,i',j',n,t} &\le x_{i,j,n,t}, \\
  \zeta_{i,j,i',j',n,t} &\le x_{i',j',n,t}, \\
  \zeta_{i,j,i',j',n,t} &\ge x_{i,j,n,t} + x_{i',j',n,t} - 1.
\end{align*}

% ============================================================
\section{Constraints}
% ============================================================

\subsection{C1 --- Placement Integrity}
\label{c1}

\begin{align}
  \sum_{n\in\cN} x_{i,j,n,t} &= a_i
  \quad \forall\,i\in\cT,\;j\in\cW_i,\;t\in\cH.
  \tag{C1}
\end{align}

Each workload is placed on \emph{exactly} one node if and only if tenant
$i$ is admitted.  The equality ($=$) is critical: it forces placement upon
admission, preventing the degenerate solution in which a tenant earns
revenue with zero infrastructure use.

\begin{impl}
  The original code used $\le$ instead of $=$, allowing the solver to
  admit all tenants without placing any workloads.  With $\le$, $\sigma$
  was unconstrained at $1.0$ and the objective appeared optimal at a
  trivially low value.  Changed to $=$ in the current implementation.
\end{impl}

\subsubsection*{C1b --- Tenant-on-Node Indicator}
\begin{align}
  y_{i,n,t} &\ge x_{i,j,n,t}
  \quad \forall\,i,\,j\in\cW_i,\,n,\,t.
  \tag{C1b}
\end{align}

\subsubsection*{C1c --- Compliance / Data Residency}
\begin{align}
  x_{i,j,n,t} &= 0
  \quad \forall\,n\notin N_{\text{allowed}}^{(i,j)},\;\forall\,t.
  \tag{C1c}
\end{align}
Implemented by fixing variables to zero at model-build time.

\subsection{C2 --- Probabilistic Capacity (SOCP)}

See \eqref{eq:c2} and \S\ref{sec:socp}.

\subsection{C3 --- Isolation Primitive Selection}

\subsubsection*{C3a --- One Primitive per Admitted Workload}
\begin{align}
  \sum_{k\in\cK} w_{i,j,k,t} &= a_i
  \quad \forall\,i,\,j\in\cW_i,\,t.
  \tag{C3a}
\end{align}

\subsubsection*{C3b --- Compliance Floor}
\begin{align}
  w_{i,j,k,t} &= 0
  \quad \forall\,k < k_{\min,i},\;\forall\,j,\,t.
  \tag{C3b}
\end{align}

\subsubsection*{C3c --- Co-location Compatibility}
\begin{align}
  \zeta_{i,j,i',j',n,t} + w_{i,j,k,t} + w_{i',j',k',t} &\le 2
  \tag{C3c}
\end{align}
for all $(i,j,i',j')\in\mathcal{P}$, $n$, $t$, $k$, $k'$ such that
$\rho_{k,k'} > \tau_{i,i'}$.  Forces co-located workloads from
incompatible tenants to use sufficiently isolating primitives.

\subsection{C4 --- Control-Plane Budgets}

\subsubsection*{C4a --- etcd Write Budget}
\begin{align}
  \sum_{i,j,k} \omega_{i,k}^{\text{etcd}}\,w_{i,j,k,t}
  &\le B^{\text{etcd}}
  \quad \forall\,t.
  \tag{C4a}
\end{align}

\subsubsection*{C4b --- API-Server QPS}
\begin{align}
  \sum_{i,j,k} \omega_{i,k}^{\text{qps}}\,w_{i,j,k,t}
  &\le B^{\text{qps}}
  \quad \forall\,t.
  \tag{C4b}
\end{align}

\subsubsection*{C4c --- Service Count}
\begin{align}
  \sum_{i,j,k} \omega_{i,k}^{\text{svc}}\,w_{i,j,k,t}
  &\le B^{\text{svc}}
  \quad \forall\,t.
  \tag{C4c}
\end{align}

\subsubsection*{C4d --- Migration-Induced QPS Churn}
\begin{align}
  \sum_{i,j,n} \delta^{\text{qps}}\,m_{i,j,n,t}^{\text{mig}}
  &\le B^{\text{qps,churn}}
  \quad \forall\,t.
  \tag{C4d}
\end{align}

\subsection{C5 --- SLA Latency Satisfaction}
\label{c5}

\begin{align}
  l^0_{i,j}
  + \sum_{\substack{(i',j')\ne(i,j)}} \alpha_{i,j,i',j'}^{\text{lat}}
    \sum_n \zeta_{i,j,i',j',n,t}
  + \sum_k \beta_{i,j,k}^{\text{lat}}\,w_{i,j,k,t}
  - e_{i,j,t}
  &\le L_{i,q(i,j)} + \BigM(1-a_i)
  \tag{C5}
\end{align}
for all $i\in\cT$, $j\in\cW_i$, $t\in\cH$.  The Big-$\BigM$ term
deactivates the constraint for rejected tenants ($a_i=0$).  $e_{i,j,t}\ge 0$
absorbs excess latency and is penalised in the objective.

\begin{note}
  $\BigM = 10^4$~ms in the current implementation (maximum feasible latency
  value).  The \emph{real-time} model also uses the label ``C5'' for the
  tenant-node access constraint; these are distinct constraints in separate
  models.
\end{note}

\subsection{C6 --- Migration Linking and Disruption Budget}

\subsubsection*{C6a --- No Migrations at $t=0$}
\begin{align}
  m_{i,j,n,0}^{\text{mig}} &= 0
  \quad \forall\,i,\,j,\,n.
  \tag{C6a}
\end{align}

\subsubsection*{C6b --- Migration Triggered by Relocation}
For $t\ge 1$:
\begin{align}
  m_{i,j,n,t}^{\text{mig}}
  &\ge x_{i,j,n,t} - x_{i,j,n,t-1}
    - \Bigl(1 - \sum_{n'\ne n} x_{i,j,n',t-1}\Bigr)
  \quad \forall\,i,\,j,\,n.
  \tag{C6b}
\end{align}

\subsubsection*{C6c --- Per-Tenant Disruption Budget}
\begin{align}
  \sum_{j,n} m_{i,j,n,t}^{\text{mig}}
  &\le \Delta_i
  \quad \forall\,i\in\cT,\;t\in\cH.
  \tag{C6c}
\end{align}

\subsection{C7 --- DRF Fairness (Max-Min)}
\label{sec:c7}

\begin{align}
  \sigma &\le s_i + (1 - a_i)
  \quad \forall\,i\in\cT,\;t\in\cH,
  \tag{C7}
\end{align}
where the DRF satisfaction fraction is
\[
  s_i \;=\;
  \frac{\displaystyle\sum_{j\in\cW_i}\sum_{n\in\cN}
        d_{i,j,r^*_i}\,x_{i,j,n,t}}{Q_{i,r^*_i}},
  \qquad
  r^*_i = \arg\max_{r\in\cR}\frac{\sum_j d_{i,j,r}}{Q_{i,r}}.
\]
The term $(1-a_i)$ releases rejected tenants from dragging $\sigma$ down.
$\sigma$ is maximised in the objective (DRF max-min fairness).

% ============================================================
\section{Optimization Objective}
% ============================================================

The model minimizes the following multi-term weighted objective:

\begin{align}
  \min \quad
  &\lambda_0 \underbrace{\sum_{n,t} \pi_n\,z_{n,t}}_{\text{infra cost}}
  +\lambda_1 \underbrace{\sum_{i,j,t} p_{i,j}\,e_{i,j,t}}_{\text{SLA penalty}}
  -\lambda_2 \underbrace{\sum_i (\bar{\pi}_i - v_i^{\text{op}})\,a_i}_{\text{admission revenue}}
  \notag\\[6pt]
  &-\lambda_3 \underbrace{\sigma}_{\text{fairness}}
  +\lambda_4 \underbrace{\sum_{i,j,k,t} c_k\,w_{i,j,k,t}}_{\text{isolation cost}}
  +\lambda_5 \underbrace{\sum_{i,j,n,t} \gamma_{i,j}\,m_{i,j,n,t}^{\text{mig}}}_{\text{migration cost}}
  \label{eq:obj}
\end{align}

\paragraph{Default weights.}
\begin{center}
\begin{tabular}{@{}llll@{}}
  \toprule
  Weight & Value & Term & Rationale \\
  \midrule
  $\lambda_0$ & 1.0  & Infrastructure cost  & Always active \\
  $\lambda_1$ & 10.0 & SLA penalty          & SLA violations are expensive \\
  $\lambda_2$ & 1.0  & Admission revenue    & Revenue incentive \\
  $\lambda_3$ & 5.0  & Fairness             & Significant equity weight \\
  $\lambda_4$ & 0.5  & Isolation cost       & Minor operational overhead \\
  $\lambda_5$ & 2.0  & Migration cost       & Moderate disruption penalty \\
  \bottomrule
\end{tabular}
\end{center}

\begin{note}
  Revenue and operational cost scale with the horizon length in the
  implementation: $\bar{\pi}_i = 3|\cH|$ and $v_i^{\text{op}} = 0.5|\cH|$.
  This ensures the objective remains balanced regardless of how many time
  slots are configured.
\end{note}

% ============================================================
\section{Model Output and Interface to Real-Time Model}
% ============================================================

The plan-ahead model's primary output consumed by the real-time model is
the per-tenant, per-slot authorised node set:
\[
  A_t^{(i)} \;=\; \bigl\{n\in\cN : \textstyle\sum_j x_{i,j,n,t} \ge 1\bigr\},
  \quad \forall\,i\in\cT,\;t\in\cH.
\]

In Python:
\begin{verbatim}
  TenantAccessSchedule: dict[(tenant_id, time_slot) -> list[node_id]]

  e.g.  { (0, 0): [2, 3],   # tenant 0 authorised on nodes 2,3 at t=0
          (0, 1): [3],       # tenant 0 restricted to node 3 at t=1
          (1, 0): [0, 1] }
\end{verbatim}

At each real-time scheduling epoch for slot $t$:
\begin{verbatim}
  active_access = filter_active_access(schedule, t)
  solve(jobs, nodes, tenant_node_access=active_access)
\end{verbatim}

The real-time model enforces access via its own constraint C5 (distinct from
the plan-ahead C5): variable upper-bound $= 0$ for any job--node pair where
the tenant is not authorised at the current slot.

% ============================================================
\section{Connection to Google Cluster-Usage Traces v3}
% ============================================================

\paragraph{Parameters derived from traces.}
\begin{center}
\begin{tabular}{@{}lll@{}}
  \toprule
  Parameter & Trace source & Field \\
  \midrule
  $\mu_{i,j,r}$  & \texttt{InstanceUsage}   & \texttt{average\_usage} \\
  $\Sigma_r$     & \texttt{InstanceUsage}   & \texttt{random\_sampled\_usage} \\
  $d_{i,j,r}$   & \texttt{InstanceEvents}  & \texttt{resource\_request} (per task, not per collection) \\
  $C_{n,r}$     & \texttt{MachineEvents}   & \texttt{capacity} \\
  $q(i,j)$      & \texttt{InstanceEvents}  & \texttt{scheduling\_class} ($3\to$ guaranteed, $0\to$ burstable) \\
  $p_{i,j}$     & \texttt{InstanceEvents}  & \texttt{priority} (higher $\to$ higher penalty) \\
  $N_{i,t}$     & \texttt{CollectionEvents}& \texttt{SUBMIT} event count per \texttt{user} per 4-hour bucket \\
  \bottomrule
\end{tabular}
\end{center}

\paragraph{Parameters requiring operator specification.}
$\pi_n$, $\bar{\pi}_i$, $v_i^{\text{op}}$ (pricing);
$L_{i,q}$ (SLA targets);
$\tau_{i,i'}$, $\rho_{k,k'}$ (interference model, from benchmarks);
$\eta_{k,r}$ (primitive overhead, from gVisor/Kata benchmarks);
$\gamma_{i,j}$, $\Delta_i$ (migration policy);
$\lambda_s$ (objective weights);
$\alpha_r^{\text{oc}}$ (overcommit policy).

\paragraph{Temporal alignment.}
Google traces use 5-minute measurement windows.  Planning slots $\cH$
correspond to coarser windows (e.g.\ 4-hour blocks).  $\mu_{i,j,r}$ and
$\Sigma_r$ per slot are estimated by aggregating measurements within each
window.  This aggregation is performed by the prediction model (separate
component).

% ============================================================
\section{Complete Model Summary}
% ============================================================

\paragraph{Decision variables.}
\begin{align*}
  &x_{i,j,n,t},\;y_{i,n,t},\;z_{n,t},\;w_{i,j,k,t},\;m_{i,j,n,t}^{\text{mig}},\;a_i
    \;\in\;\{0,1\}  \\
  &e_{i,j,t},\;\eps_{n,t}^{\text{eff}} \;\ge\; 0  \\
  &\xi_{i,j,n,k,t},\;\zeta_{i,j,i',j',n,t} \;\in\;\{0,1\} \quad\text{(McCormick auxiliaries)}\\
  &\sigma \;\in\; [0,1]
\end{align*}

\paragraph{Objective.}  See \eqref{eq:obj}.

\paragraph{Constraints.}
\begin{center}
\begin{tabular}{@{}lll@{}}
  \toprule
  Label & Name & Key property \\
  \midrule
  C1    & Placement integrity     & Equality: forces placement if admitted \\
  C1b   & Tenant-on-node indicator & Drives $\eps_{n,t}^{\text{eff}}$ \\
  C1c   & Compliance / residency  & Hard node exclusions via $N_{\text{allowed}}$ \\
  C2    & Probabilistic capacity  & SOCP / Cantelli safety margin \\
  C3a   & One primitive per workload & Exactly one $k$ selected \\
  C3b   & Compliance floor        & $k \ge k_{\min,i}$ enforced \\
  C3c   & Co-location compatibility & $\rho_{k,k'} \le \tau_{i,i'}$ for co-located pairs \\
  C4a--d & Control-plane budgets  & etcd, QPS, service count, migration churn \\
  C5    & SLA latency             & Big-$\BigM$ deactivation for rejected tenants \\
  C6a   & No migration at $t=0$   & No prior state \\
  C6b   & Migration trigger       & Relocation detection \\
  C6c   & Disruption budget       & $\le\Delta_i$ migrations per tenant per slot \\
  C7    & DRF fairness            & Max-min $\sigma$ across admitted tenants \\
  \midrule
  ---   & McCormick ($\xi$)       & $\xi = x\cdot w$ linearization \\
  ---   & McCormick ($\zeta$)     & $\zeta = x\cdot x$ linearization \\
  \bottomrule
\end{tabular}
\end{center}

\paragraph{Solver.}
Gurobi MISOCP (Mixed-Integer Second-Order Cone Program).
Default settings: \texttt{TimeLimit}~$=300$~s, \texttt{MIPGap}~$=1\%$.
Pipeline settings: \texttt{TimeLimit}~$=120$~s with \texttt{LogToConsole}
enabled for medium/high complexity samples.

\paragraph{Output.}
\texttt{TenantAccessSchedule} $\to$ passed as
\texttt{tenant\_node\_access} to the real-time model's \texttt{solve()} call.

\bigskip


\end{document}
